\chapter{S4}
\label{ch:s4}

\section{What Problem Did S4 Solve?}

S4, the Structured State Space sequence model, made state space layers competitive for
long-range deep sequence modeling \citep{gu2022s4}. Its contribution was not simply
``use an SSM.'' The basic SSM equations are old. The contribution was to combine:

\begin{itemize}
  \item HiPPO-inspired long-memory initialization;
  \item a structured parameterization of the state matrix;
  \item efficient computation of long convolution kernels;
  \item a deep residual neural-network block around the SSM layer.
\end{itemize}

The central computational goal is familiar from previous chapters. Given a discretized
SSM,
\[
  \vx_t = \bar{\mA}\vx_{t-1}+\bar{\mB}\vu_t,
  \qquad
  \vy_t = \mC\vx_t+\mD\vu_t,
\]
we want to apply the convolution kernel
\[
  \mK_i = \mC\bar{\mA}^i\bar{\mB}
\]
over very long sequences. S4 asks: how can this be expressive, stable, and fast?

\section{Why Ordinary SSMs Are Not Enough}

A generic dense state matrix $\mA\in\R^{N\times N}$ creates several problems.

\paragraph{Kernel generation is expensive.}
Computing $\mC\bar{\mA}^i\bar{\mB}$ for many lags $i$ can require repeated dense matrix
multiplications. For long sequences, this is too slow.

\paragraph{Training can be unstable.}
The eigenvalues of $\bar{\mA}$ control the memory. If they drift outside the unit circle,
the recurrence can explode. If they move too far inside, memory vanishes.

\paragraph{Random initialization is poor for long memory.}
Long-context tasks need a useful range of timescales from the start. HiPPO helps, but the
raw HiPPO matrices also need to be represented in a computationally tractable way.

\paragraph{Naive recurrence loses parallelism.}
Running the recurrence step by step is natural for inference, but it gives up the training
parallelism that makes convolutions and transformers attractive.

S4 is best understood as an answer to these four problems at once.

\section{The S4 Layer At A High Level}

An S4 layer maps an input sequence to an output sequence by:

\begin{enumerate}
  \item parameterizing a continuous-time SSM;
  \item discretizing it using a step size $\Delta$;
  \item generating a long causal convolution kernel from the discretized SSM;
  \item applying that kernel to the sequence;
  \item adding nonlinearities, channel mixing, and residual connections around it.
\end{enumerate}

At the temporal-mixing core, S4 is still an LTI system:
\[
  \vy_t = \sum_{i=0}^{t-1}\mK_i\vu_{t-i}.
\]
The novelty is how $\mK_i$ is parameterized and computed.

\section{SISO SSMs In S4}

The simplest S4 presentation uses single-input single-output SSMs. For one channel:
\[
  u_t \in \R,\qquad y_t \in \R,
\]
and
\[
  \vx_t \in \R^N.
\]
The kernel is scalar:
\[
  K_i = \mC\bar{\mA}^i\bar{\mB}.
\]

For a multi-channel neural network, S4 can use separate SISO SSMs across channels, with
pointwise projections or other mixing layers before and after the SSM operation. So one
should not picture the whole network as scalar. Rather, the SSM temporal filter itself is
often channelwise, while channel interactions are handled by surrounding linear layers.

This design choice will matter when we compare S4 with S5. S5 moves toward a more direct
multi-input multi-output SSM formulation.

\section{HiPPO Initialization In S4}

S4 begins from a HiPPO-inspired continuous-time matrix. From the previous chapter, the
purpose of HiPPO is to initialize a state that can compress long histories using
polynomial projection dynamics. This gives the SSM a good memory structure before
training.

The rough pipeline is:
\[
  \text{HiPPO continuous dynamics}
  \rightarrow
  \text{structured SSM parameterization}
  \rightarrow
  \text{discretized convolution kernel}.
\]

This is important because S4's long-memory behavior is not just a consequence of using a
large state. It comes from initializing and parameterizing the state matrix so that it
has useful modes across many timescales.

\section{Normal Matrices}

To understand S4's structure, first recall diagonalization. If
\[
  \mA = \mV\bm{\Lambda}\mV^{-1},
\]
then powers of $\mA$ are easy in the eigenbasis:
\[
  \mA^i = \mV\bm{\Lambda}^i\mV^{-1}.
\]
The problem is that an arbitrary diagonalization can be numerically poorly conditioned.
If $\mV^{-1}$ is unstable, the computation may be fragile.

A normal matrix is a complex matrix satisfying
\[
  \mA^*\mA = \mA\mA^*,
\]
where $\mA^*$ is the conjugate transpose. Normal matrices are especially nice because
they can be diagonalized by a unitary matrix:
\[
  \mA = \mV\bm{\Lambda}\mV^*.
\]
Unitary matrices preserve norms, which makes this representation much better conditioned.

Diagonal matrices are the simplest normal matrices. Symmetric real matrices and Hermitian
complex matrices are also normal. S4 exploits the fact that certain HiPPO-related
matrices can be represented as a normal part plus a low-rank correction.

\section{Normal Plus Low-Rank Structure}

S4 uses a structured state matrix often described as normal plus low-rank, or NPLR:
\[
  \mA = \mA_{\mathrm{normal}} - \mP\mQ^*,
\]
where $\mA_{\mathrm{normal}}$ is normal and $\mP\mQ^*$ is low-rank. The exact signs and
forms depend on the convention, but the idea is:

\[
  \text{HiPPO-like matrix}
  =
  \text{diagonalizable normal matrix}
  +
  \text{small low-rank correction}.
\]

This is the central S4 algebraic trick. The normal part gives efficient diagonalization.
The low-rank correction preserves the useful HiPPO memory structure. Together, they make
it possible to compute long convolution kernels efficiently without throwing away the
long-memory initialization.

On a first reading, you do not need to derive the NPLR decomposition. What matters is the
modeling role:

\begin{itemize}
  \item diagonal or normal dynamics are computationally convenient;
  \item raw HiPPO dynamics are memory-useful but not directly convenient;
  \item NPLR structure lets S4 keep both advantages.
\end{itemize}

\section{Complex Diagonal Dynamics}

After transforming into a convenient basis, S4 works with complex-valued diagonal or
almost-diagonal dynamics. A diagonal state matrix has entries
\[
  \lambda_1,\ldots,\lambda_N.
\]
The kernel contains terms like
\[
  \lambda_j^i.
\]
As Chapter~\ref{ch:linear-systems-and-filters} discussed, complex eigenvalues represent
decaying oscillatory modes:
\[
  \lambda_j = r_j e^{i\omega_j}
  \quad \Rightarrow \quad
  \lambda_j^i = r_j^i e^{i\omega_j i}.
\]

This makes S4 filters rich. A long kernel can mix many decays and oscillations. In signal
terms, the layer can learn filters with a broad range of temporal and frequency responses.

The model output remains real-valued. Complex arithmetic is a convenient internal
parameterization, not a claim that the data are complex.

\section{Generating The Convolution Kernel}

For a sequence of length $L$, S4 needs the kernel
\[
  K_0,K_1,\ldots,K_{L-1}.
\]
In a scalar SISO case,
\[
  K_i = \mC\bar{\mA}^i\bar{\mB}.
\]
If $\bar{\mA}$ is diagonal,
\[
  \bar{\mA} = \diag(\bar{\lambda}_1,\ldots,\bar{\lambda}_N),
\]
then
\[
  K_i
  =
  \sum_{j=1}^{N}
  C_j B_j \bar{\lambda}_j^i.
\]
So the kernel is a sum of exponentials. This is already much easier than powering a dense
matrix.

S4's full NPLR case is more complicated, but the same broad idea remains: exploit the
structured state matrix to compute the long kernel quickly.

\section{The Cauchy Kernel Idea}

The S4 paper emphasizes efficient computation through Cauchy kernels. The details are
technical, but the conceptual path is manageable.

Rather than compute every time-domain coefficient $K_i$ by repeated powers, one can work
with the generating function or frequency-domain representation of the SSM. For diagonal
dynamics, terms of the form
\[
  \frac{1}{\omega - \lambda_j}
\]
appear when evaluating the transfer function at frequency-like points $\omega$. A matrix
with entries of this reciprocal-difference form is called a Cauchy matrix.

So the high-level statement is:

\[
  \text{SSM transfer-function evaluation}
  \quad \leadsto \quad
  \text{structured Cauchy matrix computation}.
\]

Efficient Cauchy multiplication lets S4 generate the convolution kernel or its frequency
representation much faster than a naive dense approach. For learning S4 conceptually, it
is enough to know that this is the specialized numerical routine that makes the structured
SSM practical at long sequence lengths.

\section{Convolutional Training}

During training, S4 usually uses convolution mode. For an input sequence $u_{1:L}$, it
computes a kernel $K_{0:L-1}$ and applies
\[
  y_t = \sum_{i=0}^{t-1} K_i u_{t-i}.
\]
This can be done in parallel over the sequence, often using FFT-based convolution or
related efficient methods.

This is a major advantage over a naive RNN-style recurrence. The underlying model has a
state-space interpretation, but training does not have to process one time step after
another.

The computational pattern is:
\[
  \text{generate long kernel}
  \rightarrow
  \text{apply causal convolution}
  \rightarrow
  \text{use ordinary backpropagation}.
\]

\section{Recurrent Inference}

The same S4 layer can also be run recurrently:
\[
  \vx_t = \bar{\mA}\vx_{t-1}+\bar{\mB}u_t,
  \qquad
  y_t=\mC\vx_t+\mD u_t.
\]
This is useful for online inference or autoregressive generation. Instead of storing the
whole past sequence, the model stores the current state $\vx_t$.

This duality is one of the most attractive features of LTI SSMs:
\[
  \text{parallel convolution for training}
  \quad \text{and} \quad
  \text{recurrent state update for deployment}.
\]

Transformers have a different deployment story: causal attention can cache previous keys
and values, but memory grows with context length. An SSM recurrent state has fixed size.

\section{Discretization In S4}

S4 parameterizes continuous-time dynamics and then discretizes them. As in
Chapter~\ref{ch:state-space-models}, discretization depends on a step size $\Delta$:
\[
  \bar{\mA} = e^{\Delta\mA},
  \qquad
  \bar{\mB}
  =
  \left(\int_0^\Delta e^{\tau\mA}\,d\tau\right)\mB.
\]

The step size can be learned. This gives the layer control over effective timescales.
Different channels can learn different discretizations, allowing the network to adapt the
continuous dynamics to the sampled sequence.

For practical understanding, remember:

\[
  \Delta
  \text{ controls how continuous memory modes appear at discrete sequence steps.}
\]

\section{The S4 Block}

The SSM kernel is only the temporal core. A deep S4 model wraps this core in a residual
block. A simplified block is:
\[
  \vz
  \rightarrow
  \text{normalization}
  \rightarrow
  \text{linear projection}
  \rightarrow
  \text{S4 temporal convolution}
  \rightarrow
  \text{activation and dropout}
  \rightarrow
  \text{linear projection}
  \rightarrow
  \text{residual addition}.
\]

Exact details vary across implementations. Some versions use gating, different
normalization placements, or additional feed-forward layers. The important design is the
same as in transformers: the specialized sequence operation is one component inside a
deep residual network.

This matters because the bare SSM is linear. Nonlinear modeling power comes from stacking
blocks with activations, gates, projections, and residual paths.

\section{What S4 Learns}

S4 learns the parameters that determine the convolution kernel. Depending on the
implementation, these include:

\begin{itemize}
  \item eigenvalue-like parameters controlling decay and oscillation;
  \item input and output vectors analogous to $\mB$ and $\mC$;
  \item step sizes $\Delta$;
  \item direct terms and surrounding channel-mixing projections.
\end{itemize}

The learned kernel can be very long. But unlike a CNN kernel, it is not an arbitrary list
of free parameters. It is generated by a structured dynamical system. That structure is
the inductive bias that gives S4 its efficiency and long-memory behavior.

\section{Why S4 Was Important}

S4 showed that state space layers could compete with attention and other sequence models
on difficult long-range benchmarks. The key lesson was not just empirical. It changed the
architectural picture:

\[
  \text{recurrence}
  \neq
  \text{slow RNN training by necessity}.
\]

A carefully structured recurrent model can expose a convolutional form for training while
retaining a recurrent form for inference. This is the foundation on which S5 and later
models build.

\section{Limitations And Complexity}

S4 is powerful, but it is not simple. Its implementation involves complex-valued
parameters, special matrix structure, discretization details, and specialized kernel
computation. These details make it less straightforward than a standard convolution or
vanilla RNN.

This complexity motivated simplified variants. S5 keeps the state-space approach but
removes some of S4's harder machinery by using diagonal state matrices and a parallel
scan formulation. Mamba later changes the emphasis again by adding input-dependent
selectivity.

\section{A Reading Strategy For The S4 Paper}

When reading \citet{gu2022s4}, it helps to separate the paper into layers:

\begin{enumerate}
  \item The modeling layer: an SSM defines a long convolution kernel.
  \item The memory layer: HiPPO gives a strong long-memory initialization.
  \item The algebra layer: NPLR structure makes the matrix tractable.
  \item The numerical layer: Cauchy kernel computation makes long convolution efficient.
  \item The neural-network layer: S4 is placed inside deep residual blocks.
\end{enumerate}

Trying to understand all five layers at once is difficult. The paper becomes much easier
once you can identify which layer a given equation belongs to.

\section{What To Remember}

\begin{itemize}
  \item S4 is a structured state space layer for long sequence modeling.
  \item It uses HiPPO-inspired dynamics to initialize long memory.
  \item The SSM kernel is still $\mK_i=\mC\bar{\mA}^i\bar{\mB}$.
  \item S4 uses normal-plus-low-rank structure to keep HiPPO-like dynamics tractable.
  \item Complex modes represent decays and oscillations.
  \item Training can use convolution mode; online inference can use recurrent mode.
  \item The S4 layer is embedded inside a nonlinear residual block.
  \item S4's power comes with implementation complexity, motivating S5.
\end{itemize}

\section{Exercises}

\begin{enumerate}
  \item Explain why a generic dense $\mA$ is problematic for long-sequence SSMs.
  \item In your own words, what does NPLR structure buy us?
  \item Why are complex eigenvalues natural for temporal filters?
  \item Describe the difference between S4's convolutional training mode and recurrent
  inference mode.
  \item When reading the S4 paper, identify whether a difficult equation belongs mainly
  to memory initialization, matrix structure, numerical kernel computation, or the neural
  block design.
\end{enumerate}

\section{Prepares You For}

S5 keeps the state-space viewpoint but simplifies the parameterization and changes the
multi-channel treatment. The next chapter explains how S5 moves from S4's more elaborate
NPLR machinery toward diagonal state matrices, MIMO SSMs, and parallel scan.
